% !TeX program = pdflatex
% Compile twice with pdfLaTeX. The bibliography is included below.
\documentclass[11pt,a4paper]{article}
\usepackage[T1]{fontenc}
\usepackage{lmodern,microtype}
\usepackage[margin=1in]{geometry}
\usepackage{array,booktabs,tabularx,titlesec}
\usepackage[numbers]{natbib}
\usepackage[hidelinks]{hyperref}
\hypersetup{pdftitle={Evidence Reuse and Conformer Selection in Auto-ORCA},
  pdfauthor={CHEN GENG},pdflang={en}}

% Replace \blankfield with the relevant details when ready.
\newcommand{\blankfield}{\rule{3.5cm}{0.35pt}}
\newcommand{\StudentName}{CHEN GENG}
\newcommand{\StudentID}{\blankfield}
\newcommand{\Supervisor}{\blankfield}
\newcommand{\SubmissionDate}{\blankfield} % DD/MM/YYYY

\setlength{\parindent}{0pt}
\setlength{\parskip}{0.42em}
\setlength{\emergencystretch}{1.5em}
\setlength{\tabcolsep}{5pt}
\renewcommand{\arraystretch}{1.16}
\newcolumntype{Y}{>{\raggedright\arraybackslash}X}
\titleformat{\section}{\fontsize{12}{14}\selectfont\bfseries}{\thesection.}{0.5em}{}
\titlespacing*{\section}{0pt}{1.05em}{0.35em}
\newcommand{\subhead}[1]{\par\addvspace{0.4em}\textbf{#1}\par\nobreak}
\newcommand{\unitE}{\ensuremath{\mathrm{kJ\,mol^{-1}}}}
\newcommand{\doi}[1]{\href{https://doi.org/#1}{doi:\,#1}}
\newcommand{\archivebase}{https://github.com/az87988799/BG6022-v4/blob/9aaa4ae40dabe83ef496f71b4434f7f4987b2089}
\raggedbottom
\begin{document}

\textbf{BG6023 MSc Research I}\hfill CA1 Research Proposal\par
\vspace{0.25em}
{\fontsize{12}{15}\selectfont\bfseries
Evidence Reuse and Conformer Selection in Auto-ORCA\par}
\begin{tabularx}{\linewidth}{@{}l Y l Y@{}}
Student name & \StudentName & Student ID & \StudentID \\
Supervisor   & \Supervisor  & Date of submission & \SubmissionDate
\end{tabularx}

\section{Introduction and Background}
An ORCA output can answer several questions about a molecule, but the answer
depends on what was calculated. Two single-point energies at the same geometry
can show how an implicit solvent changes the electronic energy. A researcher
asking for a hydration free energy also needs a thermodynamic model and defined
standard states. This distinction matters when a large language model (LLM)
interprets density functional theory (DFT) calculations for a user.
ORCA supplies the quantum chemistry methods~\cite{neese}. Jacobs and Pollice
tested whether an LLM could generate executable ORCA inputs~\cite{jacobs},
and the El Agente Quntur system used ORCA results to guide the interpretation
of multistep investigations~\cite{quntur}.

Conformer screening poses a related decision. A flexible molecule can have
several low-energy shapes, and a researcher may be able to afford DFT
calculations for only a few of them. Multilevel workflows such as CENSO use
inexpensive calculations to screen candidates before refining their
energies~\cite{censo}. Auto-ORCA, the prototype developed in this project,
turns a research question into ORCA jobs, retains the input geometries and
raw outputs, and lets a single agent choose what to run next. For a laboratory
using one workstation, each additional job must justify its computational cost.

\section{Problem Statement}
In an earlier trial of Auto-ORCA, seven of eight ORCA jobs terminated normally,
yet none of the seven research requests met all its delivery requirements.
For example, six completed points in a butane torsion scan still left a
requested endpoint missing~\cite{trial}. This makes the agent's decision
after a calculation a natural subject for evaluation.

In the project's ethanol pilot, B3LYP calculations at fixed geometry gave
water-minus-gas electronic energy differences of $-16.61$ and $-18.28$~\unitE{}
with def2-SVP and def2-TZVP, respectively. Water was represented by a
conductor-like polarizable continuum model (CPCM). Those values establish
the sign of the electronic effect. A standard hydration free energy remains
undetermined from those records. In a separate pilot, ORCA ranked two sets
of four and six fixed \emph{n}-hexane geometries. The six-candidate set spanned
$12.61$~\unitE{}~\cite{archive}. Both sets were evaluated in full. The open
question is how reliably the agent can find the lowest-energy geometry in
a fresh set with just four queries.

\section{Aim, Objectives and Research Questions}
\textbf{Aim.} To test how reliably Auto-ORCA can answer questions from existing
DFT evidence and allocate a fixed number of new calculations when that evidence
is incomplete.

\textbf{Objective 1.} Measure the proportion of eight question pairs answered
correctly with and without an explicit check of the available evidence.
\emph{RQ1. Does checking the requested quantity and calculation conditions
reduce unsupported answers while preserving answers the files can support?}

\textbf{Objective 2.} Compare strategies for selecting four geometries for
DFT evaluation on six held-out molecules.
\emph{RQ2. How often does feedback from early DFT energies help Auto-ORCA
recover the lowest-energy candidate, and how much energy does it miss?}

\clearpage
\section{Methodology}
\subhead{Study 1 \enspace Answering paired questions from the same files}
I will prepare one gas-phase and one CPCM-water B3LYP/def2-SVP single-point
output for each of eight neutral, closed-shell organic molecules held out from
development. Each pair of calculations will use identical coordinates. For four
molecules the two questions will reverse the direction of the energy difference.
For two more, the target will change from that difference to a standard
hydration free energy. For the final two, I will withhold the water output
and ask first for the available gas-phase energy and separately for the solvent
effect. Each pair thus shares exactly the same files. Ethanol will serve as a
development example.

I will compare the existing agent with the same agent prompted to identify the
target quantity, match geometry and conditions, and cite the source of each
value before deciding whether further work is needed. Both will have the same
read-only files, tools, model version and call limit. Each question will be
asked in a fresh session twice per condition (64 sessions), with the order of
conditions counterbalanced. A request to rerun ORCA will be recorded without
launching a job. I will write a reference answer from each raw output and
have a second chemistry reader check the values and missing prerequisites
before either condition sees a question.

A pair passes a repeat only if both answers identify what the files support.
For a numerical answer I will require the correct sign, units, supporting
output lines and agreement within $0.1$~\unitE{} for a difference, or
$10^{-5}$~hartree for a total energy. When evidence is insufficient,
the answer must withhold the number and name the missing calculation or
thermodynamic information. I will score anonymised responses
against the frozen key and report paired counts per molecule. I will also
count unsupported numbers among unanswerable questions, and omissions and
unnecessary rerun requests among answerable questions.

\subhead{Study 2 \enspace Selecting four conformers without seeing the full ranking}
For six additional neutral, flexible molecules containing C, H and O, I will
generate candidates with a fixed embedding seed and relax them with the MMFF94
force field. I will remove duplicates using a symmetry-aware heavy-atom
root-mean-square displacement cutoff of 0.35~\AA{} and retain the first eight
distinct geometries in generation order. MMFF94 and the GFN2-xTB tight-binding
method will rank those same coordinates. I will compute all 48 gas-phase
$r^2$SCAN-3c single-point reference energies, verify convergence, and withhold
the table from the selectors. Looking up one energy will count as one DFT
evaluation for every selector.

The first rule takes the four lowest MMFF94 candidates. The second takes the
lowest member of each distinct torsion family in MMFF94 order, then fills spare
places from the same ranking. Families will be defined in advance by the
anti and two gauche regions of designated torsions. In two further conditions,
Auto-ORCA receives both rankings and the family labels and starts with the
MMFF94 minimum. In one condition it commits to the remaining three choices
immediately. In the other it sees the DFT energy after each choice. This
comparison isolates the information in the feedback. Each fixed rule runs
once. Each agent condition runs three times per molecule in fresh sessions
with the same model and limits. A session that ends before making four
selections counts as a miss. I will report its number of completed
queries and mark the four-query energy comparison as missing.

\clearpage
\subhead{How to score a conformer selection}
For each molecule I will score a hit when the best selected geometry lies
within $0.01$~\unitE{} of the lowest reference energy in its eight-candidate
set. I will also report the energy missed, in \unitE{},
\[
  \min_{c\in Q} E_{\mathrm{DFT}}(c)-\min_{c\in S} E_{\mathrm{DFT}}(c),
  \qquad |Q|=4,\quad |S|=8,
\]
where $Q$ contains the four queried geometries and $S$ all eight candidates.
I will average repeats within each molecule and report paired differences
in hit rate and energy missed, alongside individual cases. I will record
model calls, tokens and CPU core-hours, costing reference generation once.

\section{Expected Outcomes and Significance}
I expect the explicit evidence check to help most when a familiar output is
used to ask for a different physical quantity. For conformers, I expect the
adaptive procedure to help chiefly when the two inexpensive rankings disagree.
If the fixed rules find the same minima at lower cost, they would be preferable
for those molecules.
The deliverables will be a scored set of follow-up questions, six candidate
energy tables and a record of each selection. These will let computational
chemists inspect a recommendation against the underlying calculation. Findings
will apply to the tested model, molecular set and fixed geometries. Establishing
equilibrium populations would require additional sampling and free energies.

\section{Timeline}
\begin{tabularx}{\linewidth}{@{}>{\raggedright\arraybackslash}p{0.34\linewidth}
  c c Y@{}}
\toprule
\textbf{Task} & \shortstack{\textbf{Start}\\(week)} & \shortstack{\textbf{End}\\(week)} & \textbf{Output} \\
\midrule
Review literature and pilots    & 1 & 2 & Case definitions and scoring rules \\
Pilot and freeze protocols      & 3 & 4 & Checked keys, candidate lists and budget \\
Collect paired answers          & 5 & 6 & Anonymised answers and decision logs \\
Evaluate conformer selection    & 5 & 8 & DFT reference tables and query sequences \\
Audit and analyse outcomes      & 9 & 10 & Paired comparisons and cost tables \\
Write up and present            & 11 & 12 & Research paper and presentation \\
\bottomrule
\end{tabularx}

\section{Resources}
The study requires the local ORCA~6.1.1 installation, the ORCA Python
Interface, Python, RDKit, xTB, one LLM endpoint and a chemistry reader.
The workstation will run one ORCA job at a time on up to four cores with
1~GiB total memory and 192~MB MaxCore per process. The two studies require
16 solvent/gas jobs and 48 conformer reference jobs. I will cap the total
at 80 ORCA starts including pilots and retries, and each of the 100 planned
agent sessions at eight model calls. If a reference fails, I will log the cause and
use a prespecified reserve molecule only before unblinding. An incomplete
set will be reported if the budget expires. The study uses generated molecular
structures and involves no human participants. I will restrict model access
to the study files, protect credentials and share inputs and scoring scripts
under their applicable licences.

\clearpage
\renewcommand{\bibsection}{\section{References}}
\setlength{\bibsep}{0.65em}
\begin{thebibliography}{9}
\bibitem{neese}
F. Neese. Software update: The ORCA program system---version 6.0.
\emph{WIREs Computational Molecular Science} 15 (2025), e70019.
\doi{10.1002/wcms.70019}.

\bibitem{jacobs}
P. F. Jacobs and R. Pollice. Developing large language models for quantum
chemistry simulation input generation. \emph{Digital Discovery} 4 (2025),
762--775. \doi{10.1039/D4DD00366G}.

\bibitem{quntur}
J. B. P\'erez-S\'anchez, Y. Zou, J. A. Campos-Gonzalez-Angulo, et al.
El Agente Quntur: A research collaborator agent for quantum chemistry.
\emph{Chem} (2026), 103259. \doi{10.1016/j.chempr.2026.103259}.

\bibitem{censo}
S. Grimme, F. Bohle, A. Hansen, P. Pracht, S. Spicher and M. Stahn.
Efficient quantum chemical calculation of structure ensembles and free
energies for nonrigid molecules. \emph{Journal of Physical Chemistry A}
125 (2021), 4039--4054. \doi{10.1021/acs.jpca.1c00971}.

\bibitem{trial}
Auto-ORCA project. \emph{Seven-task evaluation of the frozen Auto-ORCA
workflow} (2026). Experiment report, 20--21 October.
\href{\archivebase/docs/acceptance/code-as-tools-step7-refreeze-3/report.md}
{Per-task execution records and independent reviews}.

\bibitem{archive}
Auto-ORCA project. \emph{Archived ORCA calculation results} (2026).
Project data at revision \texttt{9aaa4ae4}.
\href{\archivebase/FinalReport-LaTeX/data/T5-environments.csv}{Ethanol
solvent and gas energies (T5)} and
\href{\archivebase/FinalReport-LaTeX/data/T2-candidates.csv}{\emph{n}-hexane
candidate energies (T2)}.
\end{thebibliography}
\end{document}
